I have everything I need from the fact sheet. Here is the Discussion section.

\section{Discussion}

The design choice that most shapes how this detector behaves is that every prediction is read out from the last timestep of a strictly causal, left-padded network. A causal model at inference sees only past and present samples, which is exactly what the offline evaluation replays: a window ending at time $t$ produces the same logit whether it arrives from a recorded CSV or from the live \texttt{/lowstate} stream. That equivalence is not an assumption but a measured property. Over the test windows the runtime and research paths agree to within $3.8\times10^{-7}$ on the detection logit and $2.1\times10^{-7}$ on intensity, so the numbers reported here are the numbers the deployed node would emit on the same inputs. The practical payoff is that the ONNX node runs in roughly $1.0$~ms per window (p95 $\sim1.3$~ms) against a $2.0$~ms budget at $500$~Hz, which leaves the detection loop comfortably real-time on a single CPU thread. This latency was measured on a development CPU, not the Go2's onboard computer, so it should be read as a benchmark of the deployment runtime rather than an on-robot figure.

Temporal integration is what makes the task tractable at all, and the baseline makes the argument for us. The entanglement signature is subtle: a modest rise in resistive thigh torque against little joint motion, easily confused with a firm stance or a momentary load. An instantaneous thigh-torque rule captures the right physics but fires on single samples, and its recall collapses to $0.131$ even though its precision ($0.735$) is respectable. The TCN integrates that faint pattern over a $254$~ms receptive field, and recall rises to $0.818$ at comparable precision ($0.799$), lifting F1 from $0.222$ to $0.809$ and ROC-AUC from $0.830$ to $0.956$. The point is not that the heuristic is wrong about the mechanism, but that a single timestep does not carry enough evidence to separate a snag from normal loading; the window does.

Per-leg attribution exposes the honest cost of a small dataset. Front-right entanglement appears in only three positive recordings, and the effect is visible directly in the operating point: FR reaches perfect recall ($1.000$) but its precision falls to $0.467$, the weakest of the four legs, whereas the better-represented FL attains $0.731$ precision at the same full recall. This is a data-coverage artifact, not a modeling failure. The model has seen too few distinct FR events to sharpen the decision boundary, so it over-flags to avoid missing the class. We would expect the trade-off to relax with a few more FR recordings rather than with any change to the architecture.

Two of the more instructive results concern how failures were fixed and how confidence should be read. Specific false-positive modes, Lock stances mistaken for snags and firing during backward walking, were removed by adding targeted hard-negative recordings: Lock false-fires dropped from $4.7\%$ to $0$ and backward-walking firing from $9.7\%$ to $0$, while the front-right event's correct-leg firing rose from $4\%$ to $100\%$, all without touching the network. That the cure was data, not depth, is encouraging for maintainability, but it is an admission too: the model had learned shortcuts that only new negatives could unlearn, and other unseen stance modes may hide similar failures. Calibration carries a matching caveat. Temperature scaling ($T=5.918$) de-saturates probabilities that were otherwise pinned near $1.0$ and slightly improves the Brier score ($0.128\rightarrow0.118$), which is what lets the operating threshold move off the $0.9999$ clamp to $0.963$. It does not improve reliability in the strict sense; ECE actually worsens ($0.134\rightarrow0.156$). The scaled confidence should therefore be treated as a usable ordering of alarm strength for thresholding and per-leg gating, not as a faithful posterior probability.